\selectlanguage{russian}
\section{Метод}
\label{sec:Method} \index{Method}

Неизвестное соответствие внутри каждой группы естественно представить скрытой
перестановкой. Тогда восстановление зависимости сводится к совместному оцениванию
условного распределения наблюдений и распределения скрытой перестановки. В работе эта идея
реализуется двумя разными способами. В нейросетевом подходе оба распределения задаются
нейронными сетями, образующими вариационный автоэнкодер, и обучаются стохастическим
градиентным методом. В подходе с явными смесями условный закон и ответственности имеют
заданную аналитическую форму, поэтому их параметры обновляются модифицированным
EM-алгоритмом без обратного распространения.

\subsection{Нейросетевой вариационный подход}

В первом подходе условное распределение $p_\theta(y\mid x)$ параметризуется нейросетевым
декодером, а вариационное распределение перестановки
$q_\phi(\sigma\mid x^n,y^n)$ --- нейросетевым энкодером вместе с
Gumbel--Sinkhorn-преобразованием. Параметры $\theta$ и $\phi$ оптимизируются совместно как
параметры вариационного автоэнкодера.

\subsubsection{Вариационная целевая функция}
\label{ssec:nelbo}

Для группы $n$ используется точное маргинальное правдоподобие со скрытой перестановкой из
формулы~\eqref{eq:global-marginal}:
\begin{equation}
p_\theta(y^n\mid x^n)
=\sum_{\sigma\in S_M}p(\sigma)
\prod_{m=1}^{M}p_\theta\bigl(y^n_{\sigma(m)}\mid x^n_m\bigr),
\qquad p(\sigma)=\frac{1}{M!}.
\label{eq:neural-exact-marginal}
\end{equation}
Прямая максимизация требует суммирования по $M!$ перестановкам. Чтобы избежать этого
перебора, вводится вариационное распределение
$q_\phi(\sigma\mid x^n,y^n)$ на всём множестве $S_M$. После замены меры и применения
неравенства Йенсена получаем
\begin{align}
\log p_\theta(y^n\mid x^n)
&=\log\mathbb{E}_{q_\phi}
\left[
\frac{p(\sigma)
\prod_{m=1}^{M}p_\theta(y^n_{\sigma(m)}\mid x^n_m)}
{q_\phi(\sigma\mid x^n,y^n)}
\right] \nonumber\\
&\geq
\mathbb{E}_{q_\phi}
\left[\sum_{m=1}^{M}\log p_\theta
\bigl(y^n_{\sigma(m)}\mid x^n_m\bigr)\right]
-\mathrm{KL}\left(
q_\phi(\sigma\mid x^n,y^n)\,\|\,\mathcal{U}(S_M)
\right)
=:\mathcal{E}^n(\theta,\phi).
\label{eq:permutation-elbo}
\end{align}
Отрицательная нижняя оценка обоснованности для всей выборки равна
\begin{equation}
\boxed{
\mathcal{L}_{\mathrm{NELBO}}(\theta,\phi)
=-\frac{1}{N}\sum_{n=1}^{N}\mathcal{E}^n(\theta,\phi).
}
\label{eq:nelbo}
\end{equation}
Разрыв между логарифмом правдоподобия и ELBO имеет точный вид
\begin{equation}
\log p_\theta(y^n\mid x^n)
=\mathcal{E}^n(\theta,\phi)
+\mathrm{KL}\left(
q_\phi(\sigma\mid x^n,y^n)\,\|\,
p_\theta(\sigma\mid x^n,y^n)
\right).
\label{eq:permutation-elbo-gap}
\end{equation}
Таким образом, в теоретической постановке скрытая переменная остаётся одной перестановкой
всей группы: перехода к независимым локальным назначениям нет. При совпадении
вариационного и истинного апостериорного распределений граница становится точной.

\subsubsection{Gumbel--Sinkhorn-параметризация перестановки}
\label{ssec:relax}

Энкодер должен задавать распределение на перестановках, однако произвольная модель выдаёт
неограниченную матрицу оценок. Перестановочная матрица
$P\in\{0,1\}^{M\times M}$ содержит одну единицу в каждой строке и каждом столбце. Её
непрерывной релаксацией служит многогранник Биркгофа
\begin{equation}
\mathcal{B}_M=\left\{
P\in\mathbb{R}_{\geq 0}^{M\times M}:
\sum_{k=1}^{M}P_{mk}=1\ \forall m,\quad
\sum_{m=1}^{M}P_{mk}=1\ \forall k
\right\}.
\label{eq:birkhoff-polytope}
\end{equation}
По теореме Биркгофа--фон Неймана $\mathcal{B}_M$ является выпуклой оболочкой матриц
перестановок, поэтому его элементы интерпретируются как мягкие биективные сопоставления.

Пусть $\Phi=(\ell_{mk})\in\mathbb{R}^{M\times M}$ --- матрица логитов энкодера. Для
дифференцируемой выборки к её элементам добавляется независимый шум
$G_{mk}\sim\operatorname{Gumbel}(0,1)$, после чего применяется Sinkhorn-нормализация
~\cite{mena2018sinkhorn}. В логарифмическом пространстве одна итерация имеет вид
\begin{align}
\log\alpha_{mk}&\leftarrow\log\alpha_{mk}
-\operatorname{logsumexp}_{k'}(\log\alpha_{mk'}),\\
\log\alpha_{mk}&\leftarrow\log\alpha_{mk}
-\operatorname{logsumexp}_{m'}(\log\alpha_{m'k}).
\end{align}
Повторение нормализаций строк и столбцов приближает экспоненту результата к
$\mathcal{B}_M$. В режиме обучения и при детерминированном выводе используются
соответственно
\begin{equation}
z=\operatorname{Sinkhorn}\!\left(\frac{\Phi+G}{\tau}\right),
\qquad
z^{\mathrm{det}}=\operatorname{Sinkhorn}\!\left(\frac{\Phi}{\tau}\right).
\label{eq:stochastic-deterministic-z}
\end{equation}
При конечной температуре $\tau>0$ матрицы являются мягкими перестановками из
$\mathcal{B}_M$, а при $\tau\to0$ концентрируются около вершин, соответствующих жёстким
перестановкам. Слишком малая температура ухудшает численную устойчивость матричного
масштабирования~\cite{schmitzer2019stabilized}. Одна Sinkhorn-итерация требует $O(M^2)$
операций, а $T$ итераций --- $O(TM^2)$.

\begin{figure}[H]
\centering
\includegraphics[width=0.95\linewidth]{Latent_transform.png}
\caption{Преобразование читается слева направо: к попарным логитам энкодера добавляется
шум Gumbel, после чего Sinkhorn-нормализация строк и столбцов формирует мягкую матрицу
перестановки. При детерминированном выводе шум не добавляется.}
\label{fig:latent-transform}
\end{figure}

\subsubsection{Нейросетевой энкодер}
\label{ssec:encoder}

Для каждой пары $(x^n_m,y^n_k)$ строится периодическое представление
\begin{equation}
\left[
\sin(2\pi x^n_m),\ \cos(2\pi x^n_m),\
\sin(2\pi y^n_k),\ \cos(2\pi y^n_k)
\right]\in\mathbb{R}^4.
\end{equation}
Общая для всех пар полносвязная сеть преобразует это представление в логит
$\ell^n_{mk}$. Все логиты образуют матрицу $\Phi^n$, которая затем обрабатывается
Gumbel--Sinkhorn-преобразованием~\eqref{eq:stochastic-deterministic-z}. Параметры
$\phi$ общие для всех групп, поэтому энкодер амортизирует вывод: для новой группы не
требуется отдельная оптимизация вариационных параметров.

\begin{figure}[H]
\centering
\includegraphics[width=0.98\linewidth]{Encoder.png}
\caption{Нейросетевая параметризация попарных логитов вариационного распределения.}
\label{fig:encoder-architecture}
\end{figure}

\subsubsection{Нейросетевой декодер}
\label{ssec:decoder}

Условное распределение задаётся нейросетевой смесью из $K$ компонент:
\begin{equation}
p_\theta(y\mid x)
=\sum_{r=1}^{K}\pi_r(x)\,
\tilde{\mathcal{N}}\bigl(y\mid\mu_r(x),\epsilon^2\bigr).
\label{eq:decoder-mixture}
\end{equation}
Сеть получает периодические признаки $[\sin(2\pi x),\cos(2\pi x)]$ и возвращает средние
$\mu_r(x)$ и логиты весов $\alpha_r(x)$, причём
\begin{equation}
\pi_r(x)=\frac{\exp(\alpha_r(x))}{\sum_{s=1}^{K}\exp(\alpha_s(x))}.
\label{eq:decoder-mixture-weights}
\end{equation}
Дисперсия $\epsilon^2$ фиксируется и совпадает с уровнем шума генератора в синтетическом
эксперименте.

\begin{figure}[H]
\centering
\includegraphics[width=0.82\linewidth]{Decoder.png}
\caption{Нейросетевая параметризация средних и весов условной смеси.}
\label{fig:decoder-architecture}
\end{figure}

Данные рассматриваются на одномерном торе $\mathbb{S}^1=\mathbb{R}/\mathbb{Z}$; для
$y,\mu\in\mathbb{S}^1$ круговой остаток определяется как
$r(y,\mu)=((y-\mu+\tfrac12)\bmod 1)-\tfrac12$. Логарифм обёрнутой нормальной плотности
приближается вкладом ближайшего периодического образа:
\begin{equation}
\log\tilde{\mathcal{N}}(y\mid\mu,\epsilon^2)
\approx-\frac12\left[
\log(2\pi\epsilon^2)+\frac{r(y,\mu)^2}{\epsilon^2}
\right].
\label{eq:wrapped-gaussian-approx}
\end{equation}
При малом $\epsilon$ вклады дальних периодических образов экспоненциально малы.

Для каждой пары точек определяется модельная плотность
\begin{equation}
p_\theta(y^n_j\mid x^n_m)
=\sum_{r=1}^{K}\pi_r(x^n_m)
\tilde{\mathcal{N}}\bigl(y^n_j\mid\mu_r(x^n_m),\epsilon^2\bigr)
\label{eq:pairwise-decoder-density}
\end{equation}
и попарное отрицательное логарифмическое правдоподобие
\begin{equation}
D^n_{mj}=-\log p_\theta(y^n_j\mid x^n_m).
\label{eq:pairwise-cost}
\end{equation}
Мягкая перестановка $z^n$ задаёт реконструкционный член
\begin{equation}
\mathcal{L}_{\mathrm{rec}}^n
=\sum_{m=1}^{M}\sum_{j=1}^{M}z^n_{mj}D^n_{mj}.
\label{eq:recon}
\end{equation}

\subsubsection{Обучение стохастическим градиентным методом}
\label{ssec:sgd}

Формулы~\eqref{eq:permutation-elbo}--\eqref{eq:nelbo} задают точную вариационную цель на
пространстве перестановок. В реализации распределение Gumbel--Sinkhorn имеет неявную
плотность, поэтому точный KL на $S_M$ недоступен. Для одной группы оптимизируется
дифференцируемый суррогат
\begin{equation}
\mathcal{L}^n(\theta,\phi)
=\underbrace{\sum_{m,k}z^n_{mk}D^n_{mk}}_{\text{реконструкция}}
+\beta\underbrace{\frac{1}{M}\sum_{m,k}z^{n,\mathrm{det}}_{mk}
\bigl(\log z^{n,\mathrm{det}}_{mk}+\log M\bigr)}_{\text{построчный KL-суррогат}}.
\label{eq:per-group-loss}
\end{equation}
Реконструкционный член оценивается одной Gumbel--Sinkhorn-выборкой, а KL-регуляризатор
вычисляется по строкам детерминированной мягкой перестановки. Это практическая непрерывная
аппроксимация оптимизации точной перестановочной NELBO, но не замена самой теоретической
скрытой переменной набором независимых назначений.

Обе матрицы из~\eqref{eq:stochastic-deterministic-z} дифференцируемы по $\Phi$. В основной
реализации прямой проход выполняется фиксированным числом Sinkhorn-итераций, а обратный
проход --- неявным дифференцированием неподвижной точки~\cite{eisenberger2022implicit}.
Разворачивание итераций средствами автоматического дифференцирования сохранено как
альтернативный вариант. Градиенты средней потери мини-батча ограничиваются по норме и
используются оптимизатором Adam. Вес $\beta$ линейно увеличивается от $0$ до $1$, тогда как
температура $\tau$ фиксирована на протяжении всего запуска.

После обучения новая группа обрабатывается одним прямым проходом энкодера и
детерминированной Sinkhorn-проекцией. Для извлечения жёсткой перестановки может применяться
Hungarian matching; используемый в диагностических экспериментах построчный $\arg\max$ не
гарантирует биективность.

\begin{algorithm}[H]
\caption{Нейросетевой вариационный подход с Gumbel--Sinkhorn-релаксацией}
\label{alg:catvae}
\begin{algorithmic}[1]
\Require группы $\{(x^n,y^n)\}_{n=1}^{N}$; расписание $\beta(t)$; фиксированная
температура $\tau$; фиксированная дисперсия $\epsilon^2$
\State инициализировать параметры энкодера $\phi$ и декодера $\theta$
\For{каждый шаг обучения $t$}
  \State выбрать мини-батч $\mathcal{B}$ групп
  \For{каждая группа $n\in\mathcal{B}$}
    \State $\Phi^n\gets\operatorname{Encoder}_\phi(x^n,y^n)$
    \State $z^n\gets\operatorname{Sinkhorn}((\Phi^n+G)/\tau)$
    \State $z^{n,\mathrm{det}}\gets\operatorname{Sinkhorn}(\Phi^n/\tau)$
    \State вычислить $D^n_{mk}=-\log p_\theta(y^n_k\mid x^n_m)$
    \State вычислить $\mathcal{L}^n$ по~\eqref{eq:per-group-loss}
  \EndFor
  \State обновить $(\theta,\phi)$ градиентным шагом по средней потере мини-батча
\EndFor
\State \textbf{вывод:} применить детерминированную Sinkhorn-проекцию и извлечь перестановку
\end{algorithmic}
\end{algorithm}

\subsection{Подход с явными смесями}
\label{ssec:em}

Во втором подходе нейросетевые энкодер и декодер не используются. Условное распределение
задаётся глобальной смесью сдвигов, а распределение скрытых переменных --- явными
ответственностями назначения и компоненты. Благодаря этой параметризации параметры можно
обновлять аналитически в модифицированном EM-алгоритме, не применяя SGD и обратное
распространение.

\subsubsection{Явная параметризация распределений}

Для каждой исходной точки вводится номер компоненты
$c_m\sim\operatorname{Cat}(\pi_1,\dots,\pi_K)$, после чего
\begin{equation}
y_{\sigma(m)}\sim
\tilde{\mathcal{N}}\bigl(x_m+s_{c_m},\epsilon^2\bigr).
\label{eq:shift-model}
\end{equation}
Сдвиги $s_1,\dots,s_K$ и веса $\pi_1,\dots,\pi_K$ общие для всех точек и групп;
$\theta=(s,\pi)$. В отличие от нейросетевого декодера, средние имеют заданный вид
$x+s_c$, поэтому число оцениваемых параметров не зависит от сложности нейронной сети.

Для одной группы попарное логарифмическое правдоподобие и условная ответственность
компоненты равны
\begin{align}
L_{mk}&=\log\sum_{c=1}^{K}\pi_c
\tilde{\mathcal{N}}(y_k\mid x_m+s_c,\epsilon^2),\\
w_{mkc}&=
\frac{\pi_c\tilde{\mathcal{N}}(y_k\mid x_m+s_c,\epsilon^2)}
{\sum_{c'}\pi_{c'}\tilde{\mathcal{N}}(y_k\mid x_m+s_{c'},\epsilon^2)}.
\label{eq:em-component-responsibility}
\end{align}
Здесь $w_{mkc}$ вычисляется аналитически из текущих параметров и играет роль явного
факторизованного распределения компоненты при фиксированной паре $(m,k)$.

\subsubsection{Модифицированный E-шаг}

Точное апостериорное распределение назначения требует суммы по всем перестановкам:
\begin{equation}
q^{\mathrm{exact}}_{mk}
=\sum_{\sigma\in S_M:\,\sigma(m)=k}
\frac{\prod_i\exp L_{i,\sigma(i)}}
{\sum_{\sigma'\in S_M}\prod_i\exp L_{i,\sigma'(i)}}.
\end{equation}
В модифицированном EM эта часть E-шага заменяется энтропийно-регуляризованным
Sinkhorn-планом
\begin{equation}
q=\operatorname{Sinkhorn}(L/\tau)\in\mathcal{B}_M,
\qquad
q=\arg\max_{Q\in\mathcal{B}_M}
\left\{\langle Q,L\rangle+\tau H(Q)\right\}.
\label{eq:em-e}
\end{equation}
Совместная ответственность назначения и компоненты факторизуется как
\begin{equation}
r_{mkc}=q_{mk}w_{mkc}.
\label{eq:em-joint-responsibility}
\end{equation}
Таким образом, компонентная часть E-шага имеет аналитический вид, а распределение
перестановки вычисляется посредством Sinkhorn-аппроксимации вместо факториального перебора.

\subsubsection{Аналитический M-шаг}

После суммирования ответственностей по всем группам веса смеси обновляются формулой
\begin{equation}
S_c=\sum_{n,m,k}r^{(n)}_{mkc},
\qquad
\pi_c\leftarrow\frac{S_c}{\sum_{c'}S_{c'}}.
\label{eq:em-m-pi}
\end{equation}
Сдвиг обновляется круговым средним разностей $\delta_{mk}=y_k-x_m$:
\begin{equation}
s_c\leftarrow\frac{1}{2\pi}\operatorname{atan2}\!\left(
\sum_{n,m,k}r^{(n)}_{mkc}\sin(2\pi\delta_{mk}),
\sum_{n,m,k}r^{(n)}_{mkc}\cos(2\pi\delta_{mk})
\right)\bmod1.
\label{eq:em-m-shift}
\end{equation}
Оба обновления вычисляются непосредственно и не требуют градиентных шагов. При нулевой
ответственности компоненты её предыдущий сдвиг сохраняется.

Алгоритм является модифицированным EM: точное апостериорное распределение перестановки
заменено Sinkhorn-планом, а для обёрнутой нормальной плотности круговое моментное обновление
не обязано совпадать с точным максимумом правдоподобия. Поэтому стандартная гарантия
монотонного роста правдоподобия классического EM здесь не применяется. Температура $\tau$
фиксирована на протяжении каждого запуска; в основных экспериментах используется
$\tau=1$, а влияние других фиксированных значений исследуется отдельно. Сдвиги
инициализируются по модам гистограммы попарных круговых разностей, а при невозможности
выделить моды --- почти равномерно.

\begin{algorithm}[H]
\caption{Модифицированный EM для глобальной смеси сдвигов}
\label{alg:em}
\begin{algorithmic}[1]
\Require группы $\{(x^n,y^n)\}_{n=1}^{N}$; число компонент $K$; фиксированная
температура $\tau$; фиксированная дисперсия $\epsilon^2$
\State инициализировать сдвиги $s_c$ и положить $\pi_c=1/K$
\For{итерация $t=0,1,\dots$}
  \State вычислить $L_{mk}$ и аналитические ответственности $w_{mkc}$
  \State выполнить модифицированный E-шаг:
  $q=\operatorname{Sinkhorn}(L/\tau)$ и $r_{mkc}=q_{mk}w_{mkc}$
  \State выполнить аналитический M-шаг: обновить $\pi_c$ по~\eqref{eq:em-m-pi}
  и $s_c$ по~\eqref{eq:em-m-shift}
\EndFor
\State \textbf{вывод:} вычислить $q$ при той же температуре и извлечь перестановку
\end{algorithmic}
\end{algorithm}
